%% pos/ai-declaration.tex
%% USP MBA em Inteligência Artificial e Big Data — mandatory AI-usage declaration
%% (template: USPSC-DeclaracaoIA.tex). Self-contained: defines its own form
%% macros (a blank answer line, an empty checkbox, a pre-marked checkbox) so
%% this file does not depend on any change to main.tex.
\newcommand{\formline}[1]{\makebox[#1]{\hrulefill}}
\newcommand{\checkbox}{\tikz[baseline=-0.5ex]{\draw (0,0) rectangle (0.32,0.32);}}
\newcommand{\checkedbox}{\tikz[baseline=-0.5ex]{\draw (0,0) rectangle (0.32,0.32); \draw (0.03,0.03) -- (0.29,0.29); \draw (0.03,0.29) -- (0.29,0.03);}}
\newcommand{\choice}[1]{\noindent\checkbox\hspace{0.12cm}#1}
\newcommand{\choicedone}[1]{\noindent\checkedbox\hspace{0.12cm}#1}

\chapter{Declaração de uso de Inteligência Artificial} \label{app:declaracao-ia}

\begin{center}
{\bfseries\fontsize{15}{16}\selectfont
Declaração de uso de INTELIGÊNCIA\\
ARTIFICIAL, cuidados éticos e adequação à\\
LGPD no TCC\par}
\vspace{0.65cm}
{\bfseries\large MBA em Inteligência Artificial e Big Data -- ICMC/USP}
\end{center}

\vspace{0.25cm}
\textbf{Aluno(a):} Wellington Baltazar de Souza

\vspace{0.48cm}
\textbf{Título do TCC:} Algorithmic Trading Enhanced by AI: Using News, Geopolitical Events and NLP-Based Sentiment in Forex Decision-Making

\vspace{1.0cm}
Declaro que, na elaboração deste Trabalho de Conclusão de Curso:

\vspace{0.45cm}
\choice{\textbf{Não utilizei ferramentas de Inteligência Artificial Generativa como apoio à elaboração do TCC.}}

\vspace{0.35cm}
\noindent\checkedbox\hspace{0.12cm}\textbf{Utilizei ferramentas de Inteligência Artificial Generativa como apoio à elaboração do TCC,}\\
\hspace*{0.35cm}conforme indicado abaixo.


\begin{enumerate}
\item \textbf{Como utilizei IA}: Marque as opções aplicáveis:

\vspace{0.32cm}
\choicedone{Revisão, aprimoramento ou tradução de textos}
\vspace{0.24cm}
\choicedone{Busca, síntese ou organização de informações e literatura}
\vspace{0.24cm}
\choicedone{Geração, revisão ou depuração de código}
\vspace{0.24cm}
\choicedone{Análise, tratamento ou interpretação de dados}
\vspace{0.24cm}
\choicedone{Apoio à estruturação de ideias, metodologia ou solução}
\vspace{0.24cm}
\choice{Geração ou aprimoramento de imagens, gráficos ou outros conteúdos}
\vspace{0.24cm}
\noindent\checkbox\hspace{0.12cm}Outro: \formline{10.95cm}

\vspace{0.42cm}
\textbf{Ferramentas utilizadas:} Kimi (2.5, 3.0); Claude (Sonnet, Fable 5); Codex (GPT 5.2, GPT 5.3, Astra); DeepSeek (3, 4 Pro)

\vspace{0.48cm}
\formline{15.25cm}

\vspace*{0.45cm}

\item \textbf{Nível de utilização}\\
\noindent Marque a opção que melhor representa o uso de IA na elaboração do TCC:

\vspace{0.38cm}
\choice{\textbf{Nível 1 -- Apoio pontual}}
Uso em atividades específicas, como revisão de texto, tradução, esclarecimento de dúvidas ou sugestões pontuais.

\vspace{0.38cm}
\choice{\textbf{Nível 2 -- Apoio significativo}}
Uso recorrente em algumas etapas, como estruturação de conteúdo, programação, análise de dados ou desenvolvimento da solução.

\vspace{0.38cm}
\choicedone{\textbf{Nível 3 -- Apoio substancial}}
Uso relevante na geração ou desenvolvimento de conteúdos, códigos, análises ou outros elementos que integram o TCC.

\vspace{0.72cm}
\item \textbf{Declaração sobre cuidados Éticos e atendimento à LGPD}\\
\textbf{Cuidados éticos:} quais foram os procedimentos éticos adotados em pesquisas envolvendo seres humanos, conforme as Normas e Diretrizes Regulamentadoras da Pesquisa Envolvendo Seres Humanos -- Resolução C N S N$^\circ$ 466/2012, Norma Operacional 001/2013, Resolução C N S N$^\circ$ 510/2016 e Resolução C N S N$^\circ$ 674/2022 (para pesquisas nacionais). Em trabalhos nos quais pesquisas com seres humanos não foram realizadas, a seção deve ser preenchida com ``Não se aplica'' ou ``Atende à Resolução C N S N$^\circ$ 510/2016''.

\vspace{0.4cm}
\choicedone{Não se aplica}
\vspace{0.12cm}
\noindent\checkbox\hspace{0.12cm}Cuidados adotados: \formline{10.7cm}

\vspace{0.42cm}
\formline{15.25cm}

\vspace{0.60cm}
Quando o estudo envolver \textbf{coleta, uso ou análise de dados pessoais}, ainda que disponíveis publicamente, as pessoas autoras devem também indicar as medidas adotadas para proteção de dados pessoais, em conformidade com a \textbf{Lei Geral de Proteção de Dados Pessoais} -- LGPD (Lei N$^\circ$ 13.709/2018), descrevendo, quando pertinente, os cuidados adotados no tratamento desses dados, tais como procedimentos de anonimização ou pseudonimização e medidas relacionadas ao armazenamento, acesso e segurança das informações. Imagens de pessoas menores de idade ou vulneráveis devem ser suprimidas ou devidamente anonimizadas no texto.

\vspace{0.38cm}
\choicedone{Não se aplica}
\vspace{0.12cm}
\noindent\checkbox\hspace{0.12cm}Medidas adotadas: \formline{10.65cm}

\vspace{0.42cm}
\formline{15.25cm}

\vspace*{0.45cm}
\item \textbf{Declaração de responsabilidade}\\

Declaro que revisei, validei criticamente e assumo a responsabilidade pelos conteúdos produzidos com o auxílio de ferramentas de Inteligência Artificial.

\vspace{0.55cm}
Declaro, ainda, que compreendo plenamente, sou capaz de explicar e defender todo o conteúdo apresentado neste TCC, bem como que observei os princípios éticos e os requisitos de proteção de dados aplicáveis.

\vspace{1.0cm}
\textbf{Data:} \formline{0.85cm} / \formline{0.85cm} / \formline{1.5cm}

\vspace{0.55cm}
\textbf{Assinatura do(a) aluno(a):} \formline{7.9cm}
\end{enumerate}
